\section{Application}
\label{sec:application}

This section documents the dual ascent method
used to solve the planner's optimisation problem~\eqref{eq:planning_problem}
at each quarter of the simulation.
It then establishes a systems-optimisation approach for the decentralised firm layer, where a set of local
linear programmes allocate production across competing firms to maximise total income
subject to distributed capital constraints, and presents performance results on the
2022 Spanish input-output benchmark. 
It must be noted that the Spanish data is used as a proof of concept with real data, not as a counterfactual policy model. 
Spain's 2022 national accounts provide a well-documented, publicly available IO table with sufficient sectoral disaggregation $n = 65$ to constitute a realistic stress test of the planning algorithm's computational and stability properties.

\subsection{Simulation Methodology}
\label{subsec:methodology}

The simulation embeds a benevolent social planner inside a fully-specified
heterogeneous multi-agent economy.
The implementation is written in Python (orchestration, calibration,
post-processing) with a compiled Julia backend (\texttt{ModelCore.jl})
that performs all hot inner loops.
To minimize cross-language overhead, we employ a \textbf{Vectorized Neumann Expansion}
technique: during calibration, firm-level capital allocations are computed as a 
single batch matrix operation in Julia, reducing inter-process communication 
latency by $99\%$ compared to sequential firm loops.

\paragraph{Calibration and initialisation.}
The model is calibrated once from the 2022 Spanish National Accounts
input-output table ($n = 65$ sectors), yielding the technology matrix
$A$, the capital-requirement matrix $B$, the labour-coefficient vector
$\vec{l}$, and the initial capital stock $K_0$.
The subsistence vector $\vec{\gamma}$ is set to 20\% of the benchmark
consumption vector; per-household preference weights $\alpha^h$ are
drawn from a Dirichlet distribution centred on the aggregate benchmark
shares and calibrated so that household income, derived from a stochastic
distribution matrix $\Phi_{h,f}$ over $n_{\text{firms}}$ firms, clears the
initial goods market.

\paragraph{Investment pre-commitment.}
At the start of each quarter $t$, the planner first assembles a growth
signal $\vec{\hat{G}}_t$---a revealed excess-demand measure from the previous quarter's tâtonnement---and
derives a required investment vector $\Delta\vec{K}_t$ via a
third-order recursive Neumann truncation of the open Leontief capital-requirements
system. Let $\mathcal{M}(\vec{v}) = B(I-\bar{A})^{-1} \vec{v}$ be the capital-requirement
operator. The investment $\vec{I}^{(k)}$ at each stage $k \in \{1,2,3\}$ is computed as:
\begin{align}
  \vec{I}^{(1)}_t &= \mathcal{M}(\vec{\hat{G}}_t \odot \vec{C}_{t-1} + g \vec{G}_{t}) \\
  \vec{I}^{(k)}_t &= \mathcal{M}(\vec{\hat{G}}_t \odot \vec{I}^{(k-1)}_t + g \vec{I}^{(k-1)}_{t, G}) \quad \text{for } k > 1
\end{align}
where $\odot$ denotes element-wise multiplication. The total pre-commitment is $\Delta\vec{K}_t = \sum_{k=1}^3 \vec{I}^{(k)}_t$. This ensures that the investment required to produce the new capital is itself accounted for, up to the third order.
Government expenditure $\vec{G}_t$ grows exogenously at a fixed step $g_{\text{step}}
= 1\%$ per quarter throughout the horizon.

\paragraph{Planner optimisation.}
Given the effective resource bounds $\vec{K}_{\text{eff}}$ and $L_{\text{eff}}$
net of investment pre-commitment, the planner solves the convex programme
in~\eqref{eq:solver problem} via the Nesterov-accelerated dual ascent
described in Sections~\ref{sec:model}--\ref{sec:optimization}.
The solver is implemented natively in Julia as \texttt{solve\_planner\_native},
which operates entirely in log-space to enforce non-negativity of shadow
prices without explicit projection.
Warm-start multipliers from the previous quarter are passed in whenever
available, drastically reducing the iteration count after Q1.
The solver returns the optimal consumption target $\vec{C}^*$, the gross
output target $\vec{X}^*$, and the sector-level shadow price vector $\vec{\pi}^*$.

\paragraph{Decentralised firm production.}
Sectoral capital $\vec{K}$ is partitioned across $n_{\text{firms}}$ firms
via a random draw that preserves aggregate sectoral totals.
Each firm solves the local LP~\eqref{eq:firm_lp} natively on the Julia side 
using the \textbf{HiGHS Dual Simplex} solver. 
By offloading the entire sectoral firm-set as a block-angular constraint 
matrix to Julia, the simulation avoids the significant overhead of 
Python-based LP interfaces (e.g., SciPy). 
This allows for the simultaneous optimization of 250 independent firms 
within milliseconds. 
The alignment between planner targets ($\vec{X}^*$) and realised 
production ($\vec{X}_{\text{actual}}$) is maintained via a 
strict $0.1\%$ primal tolerance.

\paragraph{Income distribution and household demand.}
Firm nominal incomes $Y_f = \vec{v}_{\text{eff}}^\top \vec{X}_f$, where
$\vec{v}_{\text{eff}}$ is calibrated from the first-period shadow weights
so that aggregate income equals initial GDP, are distributed to
$n_h$ households via the distribution matrix:
$Y_h = \sum_f \Phi_{h,f} Y_f$.
Each household maximises a Linear Expenditure System (LES) utility
function with individual subsistence vector $\vec{\gamma}^h$ and preference
weight vector $\alpha^h$.
Aggregate household demand is then cleared against the
physical supply $\vec{C}_{\text{actual}} = (I{-}\bar{A})\vec{X}_{\text{actual}}
- \Delta\vec{K} - \vec{G}$.

\paragraph{Intra-quarter tâtonnement (\texttt{fast\_loop}).}
The quarter is split into three equal ``months''.
Within each month the Julia routine \texttt{fast\_loop\_native} iterates
Walrasian tâtonnement: prices are adjusted multiplicatively toward
The month loop iterates Walrasian tâtonnement: prices are adjusted multiplicatively toward
market-clearing and capped at a $50\%$ per-step maximum to prevent
instability. 
The system employs \textbf{Spectral Radius-based Signal Filtering} to ensure the stability of the feedback loop. At the start of the quarter, the spectral radius $\rho(M)$ of the capital-resource operator is estimated via power iteration. This value is used to clamp the revealed preference signals:
\begin{equation}
  \text{signal}_i = -\text{clamp}\left( \xi_i \frac{\Delta P_i}{P_i}, 0, \frac{1}{1 + \rho(M)} \right)
\end{equation}
where $\xi_i$ is the sectoral responsiveness coefficient derived from the LES demand structure. This filtering prevents the "bullwhip effect" in the investment signals by ensuring the gain of the cybernetic feedback remains below the secondary production threshold.

Simultaneously, household preference vectors undergo a
log-normal AR(1) drift:
\begin{multline}
  \ln \alpha_{h,i,\tau+1}
  = \ln \alpha_{h,i,\tau}
  + \theta_{\text{drift}}\!\left(\ln \bar{\alpha}_{h,i,\tau} - \ln \alpha_{h,i,\tau}\right) \\
  + \sigma_{\text{slow}} \varepsilon_{\tau}^{\text{slow}}
  + \sigma_{\text{fast}} \varepsilon_{\tau}^{\text{fast}},
\end{multline}
where the mean-reversion target $\bar{\alpha}_{h,i,\tau}$ tracks the
planner's current estimate via the habit-formation rule
(Section~\ref{sec:optimization}).
The tâtonnement terminates when the relative price change across all
sectors falls below $\delta_p = 0.5\%$, or after a maximum of 50
iterations.
The loop returns the end-of-quarter price level $\vec{P}_{t+1}$,
monthly consumption and price arrays, and an updated aggregate
preference estimate $\hat{\vec{\alpha}}_t$.

\paragraph{Capital accumulation and state update.}
Realised gross investment equals the residual of actual output after
netting consumption, government spending, and intermediate demand.
Each firm's capital stock is updated by the standard law of motion:
\begin{equation}
  K_{f,i,t+1} = (1 - \delta) K_{f,i,t} + I_{f,i,t},
\end{equation}
with $\delta = 0.01$ per quarter.
Investment is distributed across firms in proportion to their capital
shares, preserving the block-angular structure of firm balance sheets.
Aggregate capital is recovered as $K_{i,t+1} = \sum_f K_{f,i,t+1}$,
and an accounting identity check $\vec{K} \equiv \sum_f \vec{K}_f$
is enforced after every quarter.

\paragraph{Planner preference update.}
At the close of each quarter the planner updates its aggregate preference
estimate by incorporating the revealed preference signal from the
tâtonnement, adding an $\varepsilon$-floor of $10^{-5}$ to prevent
zero-weight sectors, and renormalising to the simplex.
The evolution of the planner's estimate and its distance from the true
(unobserved) household preferences are tracked via the $\ell^2$- and
$\ell^\infty$-norm gaps, which remain bounded throughout the 20-quarter
horizon (Section~\ref{sec:price_dynamics}).

\subsection{Planner's Problem Formulation}

Substituting the growth-augmented technology and investment
pre-commitment of Sections~\ref{sec:model}--\ref{sec:optimization},
the planner's quarterly problem reduces to:
%
\begin{align}\label{eq:solver problem}
\max_{\vec{C}\geq \vec{\gamma}};\sum_i \alpha^{agg}_i \ln (C_i - \gamma^{agg}_i) \notag \\
  \quad\text{s.t.}\quad
  M\vec{C} \leq \vec{K}_{\mathrm{eff}},\quad
  \vec{P}_{L} \cdot \vec{C} \leq L_{\mathrm{eff}},
\end{align}

%
where $M = B(I-\bar{A})^{-1}$ is the capital-requirement matrix
applied directly to consumption,
$\vec{K}_{\mathrm{eff}} = \vec{K}_t - M (\Delta\vec{K}_t + \vec{G}_{t})$  and
$L_{\mathrm{eff}} = L - \vec{P}_{L} \cdot ( \Delta\vec{K}_t + \vec{G}_{t})$
are resource bounds net of investment pre-commitment, and
$\vec{P}_{L} = (I-\bar{A})^{-\top}\vec{l}$ is the
transformed labour coefficient vector precomputed once at calibration.

This is a convex programme: the objective is strictly concave and the
feasible set is a convex polytope, so the unique global maximiser exists
whenever the feasible set is non-empty \citep{boyd2004convex}.

\subsection{Dual Formulation and Stationarity}

Rather than solving the primal problem directly, the algorithm operates on the dual.
By associating non-negative Lagrange multipliers (shadow prices) $\vec{\lambda}_K \in \mathbb{R}^n_+$ and $\lambda_L \in \mathbb{R}_+$ with the capital and labour constraints respectively, we form the Lagrangian:
%
\begin{multline} \label{eq:lagrangian effective}
  \mathcal{L}(\vec{C}, \vec{\lambda}_K, \lambda_L) = \hat{W}(\vec{C})
     - \vec{\lambda}_K^\top \bigl(M\vec{C} - \vec{K}_{\mathrm{eff}}\bigr)\\
    - \lambda_L \bigl(\vec{P}_{L} \cdot \vec{C} - L_{\mathrm{eff}}\bigr).
\end{multline}
%
Taking the derivative with respect to consumption $C_i$ and setting to zero yields the closed-form first-order condition (FOC):
%
\begin{multline} \label{eq:foc}
     \frac{\alpha^{agg}_i}{C_i - \gamma^{agg}_i} - \bigl(M^\top \vec{\lambda}_K\bigr)_i - \lambda_L P_{L, i} = 0 \\ \implies \ C_i = \gamma^{agg}_i + \frac{\alpha^{agg}_i}{\pi_i},
\end{multline}
%
where $\vec{\pi}_i = (M^\top \vec{\lambda}_K)_i + \lambda_L {P}_{L,i}$ is the implied final-demand price of good $i$.
\subsection{Multiplicative Nesterov-Accelerated Dual Ascent}

The dual function $g(\vec{\lambda}_K, \lambda_L) = \max_{\vec{C}} \mathcal{L}(\vec{C}, \vec{\lambda}_K, \lambda_L)$ is minimized over $\vec{\lambda} \geq 0$ to find the optimum shadow prices.
We employ an accelerated dual-ascent scheme using \textbf{Nesterov (FISTA) momentum} applied to the multipliers in log-space. This approach, known as \textbf{Multiplicative Dual Ascent}, ensures that shadow prices remain strictly positive without requiring explicit projection steps.

The solver maintains a state in log-space $\theta_i = \ln \lambda_i$. At each iteration $t$, we evaluate the gradients at an extrapolated point $\vec{y}^{(t)}$, where the log-gradients correspond to the normalized physical slacks:
%
\begin{equation}\label{eq:slacks}
  g_{K,j}^{(t)} = -\frac{s_{K,j}}{K_{\mathrm{eff},j}}, \quad
  g_L^{(t)} = -\frac{s_L}{L_{\mathrm{eff}}},
\end{equation}
%
where $\vec{s}_K = \vec{K}_{\mathrm{eff}} - M\vec{C}(\vec{y})$ and $s_L = L_{\mathrm{eff}} - \vec{P}_{L} \cdot \vec{C}(\vec{y})$.
The shadow prices are updated via a fixed-step gradient descent in log-space:
%
\begin{align}\label{eq:dual_update}
  \ln \lambda_{K,j}^{(t+1)} &= \ln y_{K,j}^{(t)} + \sigma^{(t)} \eta_K g_{K,j}^{(t)} \notag \\
  \ln \lambda_{L}^{(t+1)} &= \ln y_{L}^{(t)} + \sigma^{(t)} \eta_L g_L^{(t)},
\end{align}
%
where the base learning rates $\eta_K \approx 0.4$ and $\eta_L \approx 0.5$ are dynamically scaled by a factor $\sigma^{(t)}$ determined via a \textbf{backtracking line search} on the dual objective value. This ensures monotonicity in dual convergence even in the presence of high sectoral curvature.

To prevent momentum-driven oscillations near the capacity singularities, we implement an \textbf{Adaptive Gradient Restart} rule. The Nesterov momentum sequence $t_k$ is reset to unity ($t_k = 1$) whenever the extrapolation direction works against the current gradient signal:
%
\begin{equation}
  \langle \ln \vec{y}^{(t)} - \ln \vec{\lambda}^{(t)}, \vec{g}^{(t)} \rangle < 0 \implies \text{Restart}.
\end{equation}
%
Otherwise, the look-ahead point is generated via the standard FISTA recurrence:
%
\begin{equation}\label{eq:extrapolation}
  \ln y_{i}^{(t+1)} = \ln \lambda_i^{(t+1)} + \left(\frac{t_k - 1}{t_{k+1}}\right) \bigl(\ln \lambda_i^{(t+1)} - \ln \lambda_i^{(t)}\bigr).
\end{equation}
%
This formulation allows the solver to handle sectors with vastly different capital intensities unit-invariantly, as the log-space steps represent relative price adjustments (percentage changes) proportional to sectoral over-utilization.


\paragraph{Complexity per Matrix-Vector Operation (MVP).}
Each iteration of~\eqref{eq:dual_update} requires at least \textbf{two Matrix-Vector Products (MVPs)}: one forward Neumann application $M\vec{C}$ for capacity evaluation, and one adjoint application $M^\top\vec{\lambda}_K$ for shadow price valuation. The backtracking line search consumes one additional adjoint MVP per trial step to evaluate the dual objective. 
The total computational cost per quarter is measured in cumulative MVPs, where successful convergence at $10^{-5}$ tolerance is typically achieved within \textbf{400--450 MVPs}.
For sparse empiric IO tables, each MVP costs $O(k \cdot \mathrm{nnz}(A))$, where $k$ is the Neumann expansion order.
\subsection{Convergence Criteria}

The iterative loop monitors the physical constraint violations and checks for mathematical stabilization.

\paragraph{KKT Tolerance.}
The algorithm terminates successfully if both the dual complementarity and primal feasibility conditions are satisfied at a high-precision threshold (typically $10^{-4}$ or $10^{-5}$):
\begin{align}
  \text{Dual: } & \lambda_{L}^{(t)} |s_L| \leq \text{tol}_d, \quad \lambda_{K,j}^{(t)}|s_{K,j}| \leq \text{tol}_d \\
  \text{Primal: } & \frac{| \min(0, s_L) |}{L_{\text{eff}}} \leq \text{tol}_p, \quad \frac{| \min(0, s_{K,j}) |}{K_{\text{eff},j}} \leq \text{tol}_p
\end{align}
where $s_L$ and $\vec{s}_K$ are the slacks net of resource bounds. 

\subsection{Decentralised Firm Production and Household Income Distribution}

To model a distributed economy, production is not computed as a single macroeconomic aggregate. Instead, sectoral capital is partitioned across $f = 1 \dots n_{\text{firms}}$ independent firms. Each firm is endowed with private capital $\vec{K}_f$ and a proportional labour endowment $L_f$.

Due to the linear nature of the value-added objective, the classical optimisation of $\vec{v}_{\text{eff}}^\top \vec{x}_f$ exhibits non-unique optima on a flat optimal face. To mathematically emulate an \textit{Optimum Hopping} mechanism---where firms select the unique vertex combination that maximally aligns with the planner's physical output targets $\vec{X}^*$---we construct a Lexicographic Linear Programme. The tie-breaker is implemented via an $L_1$ penalty using auxiliary variables $\vec{u}_f$:

\begin{align}\label{eq:firm_lp}
\max_{\vec{x}_f \geq 0,\, \vec{u}_f \ge 0} &\quad \vec{v}_{\text{eff}}^\top \vec{x}_f - \epsilon \sum_i u_{f,i} \\
\text{s.t.} &\quad B \vec{x}_f \leq \vec{K}_f \notag\\
            &\quad u_{f,i} \geq \| x_{f,i} - (\text{share}_f \cdot X^*_i) \| \notag\\
\end{align}

Here, $\epsilon$ is chosen sufficiently small (e.g., $10^{-3}$) such that it acts strictly as a tie-breaker on the optimal face. By penalising the Manhattan distance to the proportional target, the solver algorithmically finds the unique convex combination of optimal vertices that spans $\vec{X}^*$.

The total nominal income generated by these firms, $Y_{\text{firm}} = \sum_f \vec{v}_{\text{eff}}^\top \vec{x}_f$, is then distributed heterogeneously to $h = 1 \dots n_h$ households. This distribution is governed by a stochastic ownership matrix $\Phi_{h,f}$, such that each household's income is derived from its dividend share of firm revenues: $Y_h = \sum_f \Phi_{h,f} Y_f$. Consumer demand is then computed at the individual household level using a multi-agent Linear Expenditure System (LES) and dynamically clears against the actual production achievable by the firms $\vec{C}_{\text{actual}} = (I-\bar{A})\vec{X}_{\text{actual}} - (\Delta \vec{K} + \vec{G})$ during the monthly price tâtonnement.

\subsection{Simulation Results}
\label{subsec:solver_results}

\paragraph{Price and Inflation dynamics.}
The tâtonnement loop's mean price drift averages roughly $0.90\%$ per quarter,
confirming the responsiveness of the price-adjustment process.
The geometric mean quarterly inflation rate is approximately \textbf{0.006\%}, reflecting
the system's high stability under the preconditioned guidance.
The chained CPI index is maintained via:
%
\begin{equation}\label{eq:cpi_update}
  \text{CPI}_{t+1} = (1 + i_t) \text{CPI}_t.
\end{equation}
%
where $i_t$ is the Laspeyres inflation rate realized between $t$ and $t+1$.

At $n = 65$ sectors with $250$ firms and $1,000$ households, the 
total execution time per quarter (excluding initial compilation) averages approximately \textbf{0.16 seconds}. 
This represents a \textbf{massive speedup} over preceding implementations, 
achieved primarily through vectorized data exchange and the elimination of 
redundant Python-Julia bridge calls.
The memory complexity remains strictly $O(n^2)$, ensuring scalability 
for larger IO benchmarks. The solver convergence is robust, typically requiring \textbf{160--180 iterations} (totalling \textbf{400--450 MVPs}) to satisfy the strict $10^{-5}$ KKT tolerance.

\begin{figure}[htbp]
    \centering
    \includegraphics[width=\columnwidth]{figures/hist_mvps.pdf}
    \caption{Dual Ascent throughput measured in Matrix-Vector Operations (MVPs) required per quarter to achieve $10^{-5}$ precision.}
    \label{fig:solver_mvps}
\end{figure}